%%%%%%%%%%%%%%%%%%%%%%%%%%%%%%%%%%%%%%%%%%%%%%%%%%%%%%%%%%%%%%%%%%%%%%%%%%
%
% Chapter 2
%
%%%%%%%%%%%%%%%%%%%%%%%%%%%%%%%%%%%%%%%%%%%%%%%%%%%%%%%%%%%%%%%%%%%%%%%%%%

\chapter{Mathematical Theory of Symmetry}

\label{chap:symmetry}

%%%%%%%%%%%%%%%%%%%%%%%%%%%%%%%%%%%%%%%%%%%%%%%%%%%%%%%%%%%%%%%%%%%%%%%%%%

\section{Introduction}

The notion of symmetry occupies a unique position in modern theoretical
physics. Virtually every fundamental theory is constructed by specifying
its symmetry group. Classical mechanics, electrodynamics, general
relativity, quantum mechanics, quantum field theory and the Standard
Model are all formulated in terms of invariance principles.

In spontaneous symmetry breaking the mathematical structure of the
symmetry group becomes even more important. The vacuum manifold, the
number of Goldstone bosons, the Higgs mechanism and topological defects
are all determined almost entirely by properties of Lie groups and their
actions.

The purpose of this chapter is therefore to establish the mathematical
language that will be used throughout the remainder of this monograph.

Although many of the concepts originate in abstract algebra and
differential geometry, our presentation emphasizes physical intuition
without sacrificing mathematical rigor.

%%%%%%%%%%%%%%%%%%%%%%%%%%%%%%%%%%%%%%%%%%%%%%%%%%%%%%%%%%%%%%%%%%%%%%%%%%

\section{Transformations}

A transformation is a map

\[
T:X\rightarrow X
\]

which associates to every point

\[
x\in X
\]

another point

\[
T(x).
\]

In physics,

the set

\[
X
\]

may denote

\begin{itemize}

\item Euclidean space,

\item Minkowski spacetime,

\item phase space,

\item Hilbert space,

\item field configurations,

\item vacuum manifold.

\end{itemize}

Examples include

translation,

\[
x
\rightarrow
x+a,
\]

rotation,

\[
\mathbf x
\rightarrow
R\mathbf x,
\]

reflection,

\[
x
\rightarrow
-x,
\]

and Lorentz transformation,

\[
x^\mu
\rightarrow
\Lambda^\mu_{\ \nu}
x^\nu.
\]

The physical importance of transformations lies not in the individual
maps but in the algebraic relations among them.

%%%%%%%%%%%%%%%%%%%%%%%%%%%%%%%%%%%%%%%%%%%%%%%%%%%%%%%%%%%%%%%%%%%%%%%%%%

\section{Groups}

The abstract mathematical object describing symmetry is the group.

\bigskip

\textbf{Definition.}

A group is a set

\[
G
\]

equipped with a binary operation

\[
G\times G
\rightarrow
G,
\]

satisfying

\begin{enumerate}

\item closure,

\[
ab\in G,
\]

\item associativity,

\[
(ab)c
=
a(bc),
\]

\item existence of an identity element,

\[
ea=ae=a,
\]

\item existence of inverses,

\[
aa^{-1}
=
a^{-1}a
=
e.
\]

\end{enumerate}

Groups provide the mathematical language of symmetry because successive
symmetry transformations again constitute symmetry transformations.

%%%%%%%%%%%%%%%%%%%%%%%%%%%%%%%%%%%%%%%%%%%%%%%%%%%%%%%%%%%%%%%%%%%%%%%%%%

\subsection{Examples}

The integers under addition,

\[
(\mathbb Z,+),
\]

form an infinite Abelian group.

The real numbers,

\[
(\mathbb R,+),
\]

also form an Abelian group.

Rotations in three-dimensional space form the non-Abelian group

\[
SO(3).
\]

Lorentz transformations form

\[
SO(1,3).
\]

Gauge transformations of electromagnetism form

\[
U(1).
\]

The electroweak theory is based upon

\[
SU(2)\times U(1).
\]

Quantum chromodynamics employs

\[
SU(3).
\]

Already one observes that nearly every modern physical theory is
specified by its symmetry group.

%%%%%%%%%%%%%%%%%%%%%%%%%%%%%%%%%%%%%%%%%%%%%%%%%%%%%%%%%%%%%%%%%%%%%%%%%%

\section{Abelian and Non-Abelian Groups}

Two elements

\[
a,b\in G
\]

need not commute.

If

\[
ab
=
ba
\]

for every pair,

the group is called Abelian.

Otherwise,

the group is non-Abelian.

Electromagnetism is governed by an Abelian gauge group,

whereas both the weak and strong interactions possess non-Abelian gauge
symmetries.

This distinction has profound consequences.

For example,

non-Abelian gauge bosons interact with one another,

whereas photons do not at tree level.

%%%%%%%%%%%%%%%%%%%%%%%%%%%%%%%%%%%%%%%%%%%%%%%%%%%%%%%%%%%%%%%%%%%%%%%%%%

\section{Subgroups}

A subset

\[
H
\subset
G
\]

is called a subgroup if it forms a group under the same operation.

We write

\[
H<G.
\]

Examples include

\[
SO(2)
<
SO(3),
\]

and

\[
U(1)
<
SU(2).
\]

Subgroups play a central role in spontaneous symmetry breaking because
the vacuum preserves only a subgroup of the original symmetry.

One therefore writes

\[
G
\rightarrow
H.
\]

This notation should never be interpreted as a change in the underlying
Lagrangian.

Instead,

it indicates that

\[
H
\]

is the isotropy subgroup of the vacuum.

%%%%%%%%%%%%%%%%%%%%%%%%%%%%%%%%%%%%%%%%%%%%%%%%%%%%%%%%%%%%%%%%%%%%%%%%%%

\section{Invariant Subgroups}

A subgroup

\[
N
\]

is called normal if

\[
gNg^{-1}
=
N
\]

for every

\[
g\in G.
\]

Normal subgroups permit construction of quotient groups,

\[
G/N.
\]

Although quotient groups appear throughout gauge theory, spontaneous
symmetry breaking more commonly involves quotient manifolds,

\[
G/H,
\]

where

\[
H
\]

need not be normal.

The distinction between quotient groups and homogeneous spaces will be
discussed later.

%%%%%%%%%%%%%%%%%%%%%%%%%%%%%%%%%%%%%%%%%%%%%%%%%%%%%%%%%%%%%%%%%%%%%%%%%%
%%%%%%%%%%%%%%%%%%%%%%%%%%%%%%%%%%%%%%%%%%%%%%%%%%%%%%%%%%%%%%%%%%%%%%%%%%
\section{Matrix Groups}
\label{sec:matrix_groups}

Most symmetry groups encountered in physics are realized as groups of
invertible matrices.

Let

\[
M_n(\mathbb{F})
\]

denote the set of all

\[
n\times n
\]

matrices over the field

\[
\mathbb{F},
\]

where

\[
\mathbb{F}
=
\mathbb{R}
\quad\text{or}\quad
\mathbb{C}.
\]

Matrix multiplication endows many important subsets of
\(M_n(\mathbb F)\) with natural group structures.

%%%%%%%%%%%%%%%%%%%%%%%%%%%%%%%%%%%%%%%%%%%%%%%%%%%%%%%%%%%%

\subsection{General Linear Group}

\textbf{Definition.}

The general linear group is

\[
GL(n,\mathbb F)
=
\{
A\in M_n(\mathbb F)
\;|\;
\det A\neq0
\}.
\]

The group operation is ordinary matrix multiplication.

The inverse exists because

\[
\det A\neq0.
\]

The identity element is

\[
I_n.
\]

The dimension of

\[
GL(n,\mathbb R)
\]

is

\[
n^2,
\]

since every matrix element may vary independently.

%%%%%%%%%%%%%%%%%%%%%%%%%%%%%%%%%%%%%%%%%%%%%%%%%%%%%%%%%%%%

\begin{example}

For

\[
n=2,
\]

\[
A
=
\begin{pmatrix}
a&b\\
c&d
\end{pmatrix},
\]

belongs to

\[
GL(2,\mathbb R)
\]

whenever

\[
ad-bc\neq0.
\]

\end{example}

%%%%%%%%%%%%%%%%%%%%%%%%%%%%%%%%%%%%%%%%%%%%%%%%%%%%%%%%%%%%

\begin{remark}

Every finite-dimensional continuous symmetry may be represented as a
subgroup of some

\[
GL(n,\mathbb C).
\]

This remarkable result follows from Ado's theorem for Lie algebras.

\end{remark}

%%%%%%%%%%%%%%%%%%%%%%%%%%%%%%%%%%%%%%%%%%%%%%%%%%%%%%%%%%%%%%%%%%%%%%%%%%
\section{Orthogonal Groups}

Euclidean geometry is characterized by preservation of distances.

Suppose

\[
x,y\in\mathbb R^n.
\]

Their scalar product is

\[
x\cdot y
=
x^Ty.
\]

A linear transformation

\[
x'
=
Rx
\]

preserves distances whenever

\[
(Rx)^T(Ry)
=
x^Ty.
\]

Since this must hold for arbitrary vectors,

\[
R^TR
=
I.
\]

%%%%%%%%%%%%%%%%%%%%%%%%%%%%%%%%%%%%%%%%%%%%%%%%%%%%%%%%%%%%

\textbf{Definition.}

The orthogonal group is

\[
O(n)
=
\{
R
\in
GL(n,\mathbb R)
\;|\;
R^TR=I
\}.
\]

%%%%%%%%%%%%%%%%%%%%%%%%%%%%%%%%%%%%%%%%%%%%%%%%%%%%%%%%%%%%

Taking determinants,

\[
(\det R)^2=1,
\]

therefore

\[
\det R=\pm1.
\]

Matrices with determinant

\[
+1
\]

represent proper rotations,

whereas

\[
-1
\]

includes reflections.

%%%%%%%%%%%%%%%%%%%%%%%%%%%%%%%%%%%%%%%%%%%%%%%%%%%%%%%%%%%%

\subsection{Special Orthogonal Group}

The subgroup

\[
SO(n)
=
\{
R
\in
O(n)
\;|\;
\det R=1
\}
\]

contains only orientation-preserving rotations.

Its dimension is

\[
\boxed{
\dim SO(n)
=
\frac{n(n-1)}2.
}
\]

%%%%%%%%%%%%%%%%%%%%%%%%%%%%%%%%%%%%%%%%%%%%%%%%%%%%%%%%%%%%

\paragraph{Proof.}

Write

\[
R
=
I+\epsilon A,
\]

where

\[
\epsilon
\]

is infinitesimal.

Then

\[
R^TR
=
(I+\epsilon A)^T
(I+\epsilon A)
=
I
+
\epsilon(A+A^T)
+
O(\epsilon^2).
\]

Therefore

\[
A^T=-A.
\]

Thus the generators are antisymmetric matrices.

An antisymmetric

\[
n\times n
\]

matrix possesses

\[
\frac{n(n-1)}2
\]

independent entries.

Hence

\[
\dim SO(n)
=
\frac{n(n-1)}2.
\]

\hfill$\square$

%%%%%%%%%%%%%%%%%%%%%%%%%%%%%%%%%%%%%%%%%%%%%%%%%%%%%%%%%%%%%%%%%%%%%%%%%%
\section{Unitary Groups}

Quantum mechanics preserves probabilities rather than Euclidean lengths.

The relevant inner product is

\[
\langle\psi|\phi\rangle.
\]

Suppose

\[
\psi'
=
U\psi.
\]

Probability conservation requires

\[
\langle
U\psi
|
U\phi
\rangle
=
\langle
\psi
|
\phi
\rangle.
\]

Therefore

\[
U^\dagger U
=
I.
\]

%%%%%%%%%%%%%%%%%%%%%%%%%%%%%%%%%%%%%%%%%%%%%%%%%%%%%%%%%%%%

\textbf{Definition.}

The unitary group is

\[
U(n)
=
\{
U
\in
GL(n,\mathbb C)
\;|\;
U^\dagger U=I
\}.
\]

%%%%%%%%%%%%%%%%%%%%%%%%%%%%%%%%%%%%%%%%%%%%%%%%%%%%%%%%%%%%

Its Lie algebra satisfies

\[
X^\dagger
=
-X.
\]

Thus every generator is anti-Hermitian.

%%%%%%%%%%%%%%%%%%%%%%%%%%%%%%%%%%%%%%%%%%%%%%%%%%%%%%%%%%%%

The dimension is

\[
\boxed{
\dim U(n)=n^2.
}
\]

%%%%%%%%%%%%%%%%%%%%%%%%%%%%%%%%%%%%%%%%%%%%%%%%%%%%%%%%%%%%%%%%%%%%%%%%%%
\section{Special Unitary Groups}

The determinant defines a homomorphism

\[
U(n)
\rightarrow
U(1).
\]

The kernel consists of matrices satisfying

\[
\det U=1.
\]

%%%%%%%%%%%%%%%%%%%%%%%%%%%%%%%%%%%%%%%%%%%%%%%%%%%%%%%%%%%%

\textbf{Definition.}

The special unitary group is

\[
SU(n)
=
\{
U
\in
U(n)
\;|\;
\det U=1
\}.
\]

Its Lie algebra satisfies

\[
X^\dagger=-X,
\]

together with

\[
\mathrm{Tr}(X)=0.
\]

%%%%%%%%%%%%%%%%%%%%%%%%%%%%%%%%%%%%%%%%%%%%%%%%%%%%%%%%%%%%

Therefore

\[
\boxed{
\dim SU(n)
=
n^2-1.
}
\]

%%%%%%%%%%%%%%%%%%%%%%%%%%%%%%%%%%%%%%%%%%%%%%%%%%%%%%%%%%%%

Examples include

\[
SU(2),
\]

the gauge group of weak isospin,

and

\[
SU(3),
\]

the gauge group of quantum chromodynamics.

%%%%%%%%%%%%%%%%%%%%%%%%%%%%%%%%%%%%%%%%%%%%%%%%%%%%%%%%%%%%%%%%%%%%%%%%%%
\section{Examples}

\subsection{The Group \(SO(2)\)}

Every element has the form

\[
R(\theta)
=
\begin{pmatrix}
\cos\theta&
-\sin\theta\\
\sin\theta&
\cos\theta
\end{pmatrix}.
\]

Composition satisfies

\[
R(\theta_1)
R(\theta_2)
=
R(\theta_1+\theta_2).
\]

Hence

\[
SO(2)
\simeq
U(1).
\]

This group will later describe the vacuum manifold of the Abelian Higgs
model.

%%%%%%%%%%%%%%%%%%%%%%%%%%%%%%%%%%%%%%%%%%%%%%%%%%%%%%%%%%%%

\subsection{The Group \(SO(3)\)}

The rotation group of ordinary three-dimensional space possesses

\[
3
\]

independent generators,

usually denoted

\[
J_x,
\qquad
J_y,
\qquad
J_z.
\]

Their algebra is

\[
[J_i,J_j]
=
\epsilon_{ijk}
J_k.
\]

This Lie algebra underlies orbital angular momentum,
spin,
ferromagnetism,
and rotational spontaneous symmetry breaking.

%%%%%%%%%%%%%%%%%%%%%%%%%%%%%%%%%%%%%%%%%%%%%%%%%%%%%%%%%%%%

\subsection{The Group \(SU(2)\)}

Every element may be written

\[
U
=
\exp
\left(
i
\theta^a
\frac{\sigma_a}{2}
\right),
\]

where

\[
\sigma_a
\]

are the Pauli matrices.

Explicitly,

\[
\sigma_1
=
\begin{pmatrix}
0&1\\
1&0
\end{pmatrix},
\]

\[
\sigma_2
=
\begin{pmatrix}
0&-i\\
i&0
\end{pmatrix},
\]

\[
\sigma_3
=
\begin{pmatrix}
1&0\\
0&-1
\end{pmatrix}.
\]

The generators satisfy

\[
[\sigma_i,\sigma_j]
=
2i
\epsilon_{ijk}
\sigma_k.
\]

Consequently,

\[
[T_i,T_j]
=
i
\epsilon_{ijk}
T_k,
\]

where

\[
T_i
=
\frac{\sigma_i}{2}.
\]

%%%%%%%%%%%%%%%%%%%%%%%%%%%%%%%%%%%%%%%%%%%%%%%%%%%%%%%%%%%%

\begin{remark}

Although

\[
SU(2)
\]

and

\[
SO(3)
\]

possess identical Lie algebras,

their global topology differs.

Specifically,

\[
SU(2)
\simeq
S^3,
\]

whereas

\[
SO(3)
\simeq
\mathbb{RP}^3.
\]

The map

\[
SU(2)
\rightarrow
SO(3)
\]

is a two-to-one covering map.

This distinction explains the existence of spinor
representations.

%%%%%%%%%%%%%%%%%%%%%%%%%%%%%%%%%%%%%%%%%%%%%%%%%%%%%%%%%%%%%%%%%%%%%%%%%%

